% ---------------------------------------------------------------------------------------------
% Proofs for the testing results of Section \Mref{sec:testing} (sections/testing.tex).
% Notation as in sections/theory.tex and appendix_proofs.tex. References to main-text labels use \Mref.
% ---------------------------------------------------------------------------------------------
\section{Proofs for the testing results}\label{si:testing}

\paragraph{Conventions.} All levels $\alpha$ of one-sided, TOST and Holm procedures are taken below $1/2$, so that the critical constants are positive as Lemma~\ref{lem:slutsky} requires. All limits are along a sequence of pools indexed by $N$, with $H$ and $Q$ fixed; $\Phi,\varphi$ are the standard normal distribution function and density, $z_q=\Phi^{-1}(q)$, and
``$\to_p$'', ``$\Rightarrow$'' denote convergence in probability and in law. As in Supplementary Section~\SIref{si:proofs}, $a_k=\nabla\phi_k(\tau)$ is the gradient of the $k$th contrast on the scale of means, $x^{(k)}_i=a_k^\top\xi_i=Ne^{(k)}_i$, so that $|x^{(k)}_i|\le B'$ with a constant $B'$ that depends only on the bound $B$ on the derivatives of the $\phi_k$ on the compact set $K$ of Supplementary Section~\SIref{si:proofs} and on $C$, and
$\Lambda_k=a_k^\top(\hat\tau-\tau)=\sum_hW_h(\bar x^{(k)}_{s_h}-\bar x^{(k)}_h)$ with $\operatorname{Var}(\Lambda_k)=V_k$. Recall $r^2=\sum_hW_h^2/n_h$ and, from Proposition~\Mref{prop:validity}, that under (A1)--(A3) for contrast $k$,
\begin{equation}\label{eq:si-ingredients}
X_k=\frac{\hat\theta_k-\theta_k}{V_k^{1/2}}\Rightarrow N(0,1)\ \text{(under (A4))},\qquad
\varrho_k=\Big(\frac{\hat V_k}{V_k}\Big)^{1/2}\to_p1,
\end{equation}
and the second statement also holds for the floored estimator under (A5). These two facts are the only inputs to the results for a single contrast.

\subsection{A uniform Slutsky lemma}\label{si:testing-slutsky}

\begin{lemma}\label{lem:slutsky}
Let $X_N\Rightarrow N(0,1)$ and let $\varrho_N>0$ with $\varrho_N\to_p1$. For every $0<z_-<z_+<\infty$, uniformly in $z\in[z_-,z_+]$ and $\mu,a,b\in\mathbb R$,
\begin{align}
&P(X_N+\mu>z\varrho_N)=\Phi(\mu-z)+o(1),\label{eq:slut-a}\\
&P(z\varrho_N-a<X_N<b-z\varrho_N)=\big[\Phi(b-z)-\Phi(z-a)\big]_++o(1).\label{eq:slut-b}
\end{align}
\end{lemma}
\begin{proof}
Let $F_N$ be the distribution function of $X_N$. Since the limit is continuous, P\'olya's theorem \citep[Lemma~2.11]{vandervaart1998} gives $\varepsilon_N=\sup_x|F_N(x)-\Phi(x)|\to0$, and the same bound holds for the left limits $F_N(x-)$. Fix $\eta\in(0,\tfrac12)$, let $\Omega_\eta=\{|\varrho_N-1|\le\eta\}$ and $\pi_N=P(\Omega_\eta^c)\to0$. On $\Omega_\eta$ we have $z(1-\eta)\le z\varrho_N\le z(1+\eta)$.

For \eqref{eq:slut-a}: on $\Omega_\eta$, $\{X_N>z(1+\eta)-\mu\}\subseteq\{X_N+\mu>z\varrho_N\}\subseteq\{X_N>z(1-\eta)-\mu\}$, so
\[
\Phi(\mu-z(1+\eta))-\varepsilon_N-\pi_N\le P(X_N+\mu>z\varrho_N)\le\Phi(\mu-z(1-\eta))+\varepsilon_N+\pi_N,
\]
using $P(X_N>x)=1-F_N(x)$ and $1-\Phi(x)=\Phi(-x)$. The function $\Phi$ is Lipschitz with constant $\varphi(0)$, hence both bounds are within $\varphi(0)z_+\eta+\varepsilon_N+\pi_N$ of $\Phi(\mu-z)$, uniformly in $(\mu,z)$. Let $N\to\infty$ and then $\eta\to0$.

For \eqref{eq:slut-b}: on $\Omega_\eta$ the interval $(z\varrho_N-a,b-z\varrho_N)$ contains $(z(1+\eta)-a,\,b-z(1+\eta))$ and is contained in $(z(1-\eta)-a,\,b-z(1-\eta))$. For an interval $(l,u)$ with $l<u$, $P(l<X_N<u)=F_N(u-)-F_N(l)$ lies within $2\varepsilon_N$ of $\Phi(u)-\Phi(l)$, and for $u\le l$ the event is empty and $[\Phi(u)-\Phi(l)]_+=0$. The map $(s,t)\mapsto[\Phi(t)-\Phi(s)]_+$ is Lipschitz with constant $\varphi(0)$ in each argument, and replacing $z$ by $z(1\pm\eta)$ moves each endpoint by at most $z_+\eta$. The sandwich gives \eqref{eq:slut-b} as before.
\end{proof}

\subsection{Proof of Theorems~\Mref{thm:size}(a)--(c) and \Mref{thm:power}}\label{si:testing-size}

Fix $k$ and drop it from the notation. For any $\theta_0$, $\hat\theta-\theta_0=V^{1/2}(X+\mu)$ with $\mu=(\theta-\theta_0)/V^{1/2}$, hence $Z(\theta_0)=(X+\mu)/\varrho$. By \eqref{eq:si-ingredients} and Lemma~\ref{lem:slutsky}, uniformly in $\mu\in\mathbb R$ (the range of $\mu$ as $\theta_0$ varies over $\mathbb R$),
\[
P\big(Z(\theta_0)>z\big)=\Phi(\mu-z)+o(1),\qquad P\big(Z(\theta_0)<-z\big)=P(-X-\mu>z\varrho)=\Phi(-\mu-z)+o(1),
\]
the second because $-X\Rightarrow N(0,1)$ too. Adding the two tails gives (\Mref{eq:power}) of Theorem~\Mref{thm:power} for $z=z_{1-\alpha/2}$, and the two one-sided powers are the separate tails with $z=z_{1-\alpha}$. For the TOST, with $a=(\theta+\delta)/V^{1/2}$ and $b=(\delta-\theta)/V^{1/2}$ we have $Z(-\delta)=(X+a)/\varrho$ and $Z(\delta)=(X-b)/\varrho$, so the TOST rejects exactly when $z\varrho-a<X<b-z\varrho$ with $z=z_{1-\alpha}$, and \eqref{eq:slut-b} gives its power uniformly in $(a,b)$, hence in $\delta$. Consistency: if $|\mu|\to\infty$ then $\Psi_\alpha(\mu)\to1$, and the supremum bound transfers this to the rejection probability. This proves Theorem~\Mref{thm:power}.

\emph{Theorem~\Mref{thm:size}.} (a) At $\theta_0=\theta$, $\mu=0$ and the rejection probability is $2\Phi(-z_{1-\alpha/2})+o(1)=\alpha+o(1)$; moreover $Z=X/\varrho\Rightarrow N(0,1)$ by Slutsky's lemma, and applying the same computation to the level $u$ gives $P(2\{1-\Phi(|Z|)\}\le u)=P(|Z|\ge z_{1-u/2})\to u$ for every fixed $u\in(0,1)$. (b) If $\theta\le\theta_0$ then $\mu\le0$ and $P(Z>z_{1-\alpha})=\Phi(\mu-z_{1-\alpha})+o(1)\le\Phi(-z_{1-\alpha})+o(1)=\alpha+o(1)$; if $\mu\to0$ the first bound becomes $\alpha$ by continuity of $\Phi$. (c) Non-inferiority: if $\theta\ge\delta$ then $\mu'=(\theta-\delta)/V^{1/2}\ge0$ and $P(Z(\delta)<-z_{1-\alpha})=\Phi(-\mu'-z_{1-\alpha})+o(1)\le\alpha+o(1)$. TOST: if $\theta\ge\delta$ it rejects only if $Z(\delta)<-z_{1-\alpha}$, whose probability is at most $\alpha+o(1)$ by the same computation; if $\theta\le-\delta$ it rejects only if $Z(-\delta)>z_{1-\alpha}$, i.e.\ $X+a>z_{1-\alpha}\varrho$ with $a\le0$, of probability $\Phi(a-z_{1-\alpha})+o(1)\le\alpha+o(1)$. Finally let $\theta=\delta$ and $\delta/V^{1/2}\to\infty$: then $b=0$, $a=2\delta/V^{1/2}\to\infty$, and \eqref{eq:slut-b} gives $P(\text{reject})=[\Phi(-z_{1-\alpha})-\Phi(z_{1-\alpha}-a)]_++o(1)\to\alpha$; the case $\theta=-\delta$ is symmetric. $\square$

\begin{corollary}[Uniformity]\label{cor:uniform}
Let $\mathcal C_N$ be sets of pools such that every sequence $u_N\in\mathcal C_N$ satisfies (A1)--(A4) (and (A5)). Then the statements of Theorems~\Mref{thm:size}(a)--(c) and~\Mref{thm:power} hold with $\sup_{u\in\mathcal C_N}$ in front of the error terms.
\end{corollary}
\begin{proof}
If the supremum of an error term did not tend to zero, choose $u_N\in\mathcal C_N$ attaining it up to $1/N$ for every $N$. The resulting sequence satisfies the assumptions, so the error along it tends to zero, a contradiction.
\end{proof}

\subsection{Proof of Corollary~\Mref{cor:are}}\label{si:testing-are}

With $n_h=\lambda_hL\le N_h$ and $V=\sum_hN_h^2(1/n_h-1/N_h)S_h^2$ we get $V^{(d)}(L)=v_d/L-w_d$ with $v_d=\sum_hN_h^2S_h^2/\lambda_h$ and $w_d=\sum_hN_hS^2_h$ exactly. By Theorem~\Mref{thm:power} the (two-sided) power against $|\theta-\theta_0|=\Upsilon$ is $\Psi_\alpha(\Upsilon/V^{1/2})+o(1)$. The function $\Psi_\alpha$ is continuous with $\Psi_\alpha(0)=\alpha$, tends to one, and is strictly increasing on $[0,\infty)$ since $\Psi_\alpha'(\mu)=\varphi(\mu-z)-\varphi(\mu+z)>0$ for $\mu>0$, $z>0$ (because $|\mu-z|<\mu+z$). So for $\alpha<1-\beta<1$ there is a unique root $\mu_{\alpha,\beta}>0$ of $\Psi_\alpha=1-\beta$, and power $1-\beta$ holds iff $\Upsilon/V^{1/2}\ge\mu_{\alpha,\beta}$, i.e.\ $V(L)\le\Delta^2/\mu^2_{\alpha,\beta}$, i.e.\ $v_d/L\le\Delta^2/\mu^2+w_d$, which is (\Mref{eq:L-formula}). Since $\Psi_\alpha(\mu)\ge\Phi(\mu-z)$, the root satisfies $\mu_{\alpha,\beta}\le z_{1-\alpha/2}+z_{1-\beta}$; since $\Psi_\alpha(\mu)\le\Phi(\mu-z)+\Phi(-2z)$ for $\mu\ge z$, the power at $z_{1-\alpha/2}+z_{1-\beta}$ exceeds $1-\beta$ by at most $\Phi(-2z_{1-\alpha/2})$. For one-sided tests the root is exactly $z_{1-\alpha}+z_{1-\beta}$. The ratio in (\Mref{eq:are}) is algebra. For the limit, $w_dL/v_d=\sum_hN_hS^2_hL\big/\sum_hN_h^2S^2_h/\lambda_h\le\max_h\lambda_hL/N_h=\max_hf_h=:\bar f_d$, because the $h$th term of the numerator is at most $\bar f_d$ times the $h$th term of the denominator. Hence $V^{(d)}(L)=(v_d/L)(1-\vartheta_d)$ with $0\le\vartheta_d\le\bar f_d$, and $V^{(2)}(L)/V^{(1)}(L)\to v_2/v_1$ when the sampling fractions tend to zero. At the labels $L_d$ of (\Mref{eq:L-formula}) we have $v_d/L_d=\Delta^2/\mu^2+w_d$, so $w_d/(\Delta^2/\mu^2)=\vartheta_d/(1-\vartheta_d)\to0$ and the second factor in (\Mref{eq:are}) tends to one. For the Pitman efficiency: two tests with power curves $\Psi_\alpha(\Upsilon/V_1^{1/2})$ and $\Psi_\alpha(\Upsilon/V_2^{1/2})$ have the same power at every $\Upsilon$ iff $V_1=V_2$; if $V_d=v_d/L_d$ (neglecting $w_d$) then equal power at the same $\Upsilon$ needs $L_2/L_1=v_2/v_1$, the ratio of asymptotic variances (the label ratio at equal power for a fixed effect; the analogue of the Pitman efficiency of \citealt[Ch.~14]{vandervaart1998} and \citealt{noether1955}, there being no local parametric sequence). For PATTERN-DIF against SRS, $v_{\rm SRS}/v_{\rm PD}$ is the large-$N$ limit of $V^{\rm SRS}/V^{\rm PD}$, i.e.\ the efficiency of the paper. For the population target the variance is $V^{(d)}_{\rm tot}(L)=v_d/L-w_d+v_{\rm P}$ with $v_{\rm P}$ the same for all designs (Supplementary Section~\SIref{si:testing-pop}); the equation $V_{\rm tot}^{(d)}(L_d)=\Delta^2/\mu^2$ gives $L_d=\mu^2v_d/\{\Delta^2+\mu^2(w_d-v_{\rm P})\}$, and since $v_d/L-w_d=V^{(d)}(L)\ge0$ the total variance is at least $v_{\rm P}$ at every $L$, so power $1-\beta$ requires $\Delta^2\ge\mu^2v_{\rm P}$. With $w_d$ negligible, $L_2/L_1=v_2/v_1$ again. The mean squared error ratio at equal $L$ is $(V^{(2)}+v_{\rm P})/(V^{(1)}+v_{\rm P})$, which is $(\mathrm{Eff}+\varpi)/(1+\varpi)$ for $\varpi=v_{\rm P}/V^{(1)}$ and $\mathrm{Eff}=V^{(2)}/V^{(1)}$. $\square$

\subsection{A central limit lemma for stratified means}\label{si:testing-clt}

The following lemma isolates the argument of the proof of Proposition~\Mref{prop:validity}(c) for an arbitrary real variable.

\begin{lemma}\label{lem:strat-clt}
For each $N$ let $y_i$, $i\in U$, be real numbers and put $\Lambda=\sum_hW_h(\bar y_{s_h}-\bar y_h)$, $v_h=W_h^2(1-f_h)S^2_{y,h}/n_h$ and $q=\operatorname{Var}\Lambda=\sum_hv_h$. Suppose (A1) holds (so $n_h\to\infty$ and $N_h-n_h\ge\varepsilon N_h$ in strata that are not fully enumerated; fully enumerated strata have $v_h=0$), $q$ is bounded away from zero, and for every stratum either $v_h/q\to0$ or Hájek's condition holds: for every $\eta>0$,
\[
\frac{1}{(N_h-1)S^2_{y,h}}\sum_{i\in U_h:\ |y_i-\bar y_h|>\eta m_hS_{y,h}}(y_i-\bar y_h)^2\to0,\qquad m_h=\{n_h(1-f_h)\}^{1/2}.
\]
Then $\Lambda/q^{1/2}\Rightarrow N(0,1)$. The Hájek condition holds in particular if $\max_{i\in U_h}|y_i-\bar y_h|/(m_hS_{y,h})\to0$.
\end{lemma}
\begin{proof}
It suffices to show that every subsequence has a further subsequence along which the claim holds. Pass to a subsequence along which the set $\mathcal H_1$ of strata satisfying Hájek's condition and its complement $\mathcal H_2$ (where $v_h/q\to0$) are constant and $v_h/q\to\lambda_{h,\infty}^2$ for every $h$ (the ratios lie in $[0,1]$ and sum to one); then $\lambda_{h,\infty}=0$ on $\mathcal H_2$ and $\sum_{\mathcal H_1}\lambda^2_{h,\infty}=1$. The variance of $\sum_{h\in\mathcal H_2}W_h(\bar y_{s_h}-\bar y_h)/q^{1/2}$ is $\sum_{\mathcal H_2}v_h/q\to0$, so this term is $o_p(1)$. For $h\in\mathcal H_1$ put $\chi_h=(\bar y_{s_h}-\bar y_h)/\{(1-f_h)S^2_{y,h}/n_h\}^{1/2}$. By Hájek's theorem for simple random sampling without replacement \citep{hajek1960}, $\chi_h\Rightarrow N(0,1)$. The $\chi_h$ are independent across strata and $\Lambda/q^{1/2}=\sum_{\mathcal H_1}(v_h/q)^{1/2}\chi_h+o_p(1)$. The characteristic functions satisfy $\varphi_{\chi_h}(t)\to e^{-t^2/2}$ uniformly on compacts, so
$E\exp\{it\sum_{\mathcal H_1}(v_h/q)^{1/2}\chi_h\}=\prod_{\mathcal H_1}\varphi_{\chi_h}\{(v_h/q)^{1/2}t\}\to e^{-t^2/2}$. For the last claim: if $\max|y_i-\bar y_h|\le\eta m_hS_{y,h}$ for large $N$, the index set in Hájek's sum is empty.
\end{proof}

\subsection{Proof of Theorem~\Mref{thm:joint}}\label{si:testing-joint}

\emph{Linearisation.} By the argument of Supplementary Section~\SIref{si:proofs}, part (b), on the event $\mathcal E=\{\|\hat\tau-\tau\|_\infty\le c/2\}$, whose probability tends to one, $|\hat\theta_k-\theta_k-\Lambda_k|\le(B/2)C\|\hat\tau-\tau\|^2=O_p(r^2)$ for every $k\le Q$, and $V_k\ge c_kr^2$ by (A3), so $(\hat\theta_k-\theta_k-\Lambda_k)/V_k^{1/2}=O_p(r)\to_p0$. Let $\Sigma_{kl}=\operatorname{Cov}(\Lambda_k,\Lambda_l)=\sum_hW_h^2(1-f_h)n_h^{-1}\operatorname{Cov}_h(x^{(k)},x^{(l)})=\sum_hc_h\operatorname{Cov}_h(e^{(k)},e^{(l)})$; it has $\Sigma_{kk}=V_k$, and $R=D^{-1}\Sigma D^{-1}$ with $D=\operatorname{diag}(V_k^{1/2})$. Note that $V_k\le2B'^2r^2$ since $|x^{(k)}|\le B'$ (the variance bound of Supplementary Section~\SIref{si:proofs}), so $V_k\asymp r^2$.

\emph{Cram\'er--Wold.} Let $Y=D^{-1}\Lambda$ with $\Lambda=(\Lambda_1,\dots,\Lambda_Q)^\top$. The set of $Q\times Q$ correlation matrices is compact. Take a subsequence with $R_N\to R_\infty$; we show that $Y\Rightarrow N(0,R_\infty)$ along it, which proves the first claim of the theorem by the subsequence principle (every subsequence has a further subsequence with a convergent $R_N$). Fix $u\in\mathbb R^Q$ with $\|u\|=1$ and put $q_N=u^\top R_N u=\operatorname{Var}(u^\top Y)\to q_\infty=u^\top R_\infty u$. If $q_\infty=0$ then $u^\top Y\to0$ in $L^2$ and the limit $N(0,0)$ is correct. If $q_\infty>0$ write $u^\top Y=\sum_hW_h(\bar y_{s_h}-\bar y_h)$ with $y_i=b^\top x_i/r$, $x_i=(x^{(k)}_i)_k$ and $b_k=u_kr/V_k^{1/2}$; the vector $b=b_N$ is bounded because $V_k\ge c_kr^2$. Its variance is $q_N\ge q_\infty/2>0$ for large $N$, so $b^\top\Lambda_x$ has variance of exact order $r^2$, and (A4$_Q$) applies to $b_N$: for each stratum either $n_h(1-f_h)b^\top\Sigma^x_hb\to\infty$ ($\Sigma^x_h$ the stratum covariance matrix of $x$), or the stratum's share of the variance tends to zero. In the first case, $m_hS_{y,h}=\{n_h(1-f_h)b^\top\Sigma^x_hb\}^{1/2}/r$ while $\max|y_i-\bar y_h|\le2\|b\|_1B'/r$, so $\max|y_i-\bar y_h|/(m_hS_{y,h})\to0$ and Hájek's condition holds by the last sentence of Lemma~\ref{lem:strat-clt}. Lemma~\ref{lem:strat-clt} gives $u^\top Y/q_N^{1/2}\Rightarrow N(0,1)$, hence $u^\top Y\Rightarrow N(0,q_\infty)$. Since $u$ was arbitrary, the Cram\'er--Wold device gives $Y\Rightarrow N(0,R_\infty)$. With the linearisation, $\{(\hat\theta_k-\theta_k)/V_k^{1/2}\}_k=Y+o_p(1)$ has the same limit.

\emph{Covariance estimation.} Let $\hat a_k=\nabla\phi_k(\hat\tau)$ and $\hat S_h$ the sample covariance matrix of $\xi_i$ in stratum $h$. Then $\hat\Sigma_{kl}=\sum_hW_h^2(1-f_h)n_h^{-1}\hat a_k^\top\hat S_h\hat a_l$ and $|\hat\Sigma_{kl}-\Sigma_{kl}|\le r^2\max_h|\hat a_k^\top\hat S_h\hat a_l-a_k^\top\Sigma^\xi_ha_l|=r^2o_p(1)$ by Lemma~\Mref{prop:lemma-cov}, the consistency $\hat a_k-a_k\to_p0$ and boundedness on $\mathcal E$, exactly as in the proof of Proposition~\Mref{prop:validity}(d). The matrix $\hat\Sigma=\sum_hW_h^2(1-f_h)n_h^{-1}\hat A\hat S_h\hat A^\top$, with $\hat A$ the $Q\times C$ matrix of the $\hat a_k^\top$, is positive semi-definite. The floored estimator is $\hat\Sigma^\flat=\hat\Sigma+\operatorname{diag}(\hat V^\flat_k-\hat V_k)$, with non-negative increments, hence positive semi-definite, and $\hat V^\flat_k/V_k=(\hat V_k^\flat/\hat V_k)(\hat V_k/V_k)\to_p1$ by (A5) and Proposition~\Mref{prop:validity}(d). With $\hat D=\operatorname{diag}((\hat V^\flat_k)^{1/2})$ and $\hat R=\hat D^{-1}\hat\Sigma^\flat\hat D^{-1}$ (a correlation matrix),
$\hat R_{kl}-R_{kl}=\{\hat\Sigma_{kl}/(\hat V_k^\flat\hat V_l^\flat)^{1/2}\}-\Sigma_{kl}/(V_kV_l)^{1/2}\to_p0$ for $k\ne l$, because $\hat\Sigma_{kl}-\Sigma_{kl}=r^2o_p(1)$, $V_k\ge c_kr^2$ and $\hat V^\flat_k/V_k\to_p1$; and $\hat R_{kk}=R_{kk}=1$. $\square$

\subsection{The critical value of the max-\texorpdfstring{$|t|$}{t} procedure}\label{si:testing-crit}

Let $\mathcal R_Q$ be the set of $Q\times Q$ correlation matrices (symmetric, positive semi-definite, unit diagonal; compact). For $R\in\mathcal R_Q$ and $N\sim N(0,R)$ let $F_R(c)=P(\max_{k\le Q}|N_k|\le c)$. For a non-empty $S\subseteq\{1,\dots,Q\}$, $R_S\in\mathcal R_{|S|}$ is the principal submatrix.

\begin{lemma}\label{lem:crit}
For every $\alpha\in(0,1)$ and $R\in\mathcal R_Q$:
\begin{enumerate}
\item[(i)] $F_R$ is continuous and strictly increasing on $[0,\infty)$, so $c_\alpha(R)=F_R^{-1}(1-\alpha)$ is well defined;
\item[(ii)] $z_{1-\alpha/2}\le c_\alpha(R)\le z_{1-\alpha/(2Q)}$;
\item[(iii)] $c_\alpha(R_S)\le c_\alpha(R_{S'})$ for $S\subseteq S'$;
\item[(iv)] $(R,c)\mapsto F_R(c)$ is jointly continuous, and $R\mapsto c_\alpha(R)$ is continuous, hence uniformly continuous on $\mathcal R_Q$.
\end{enumerate}
\end{lemma}
\begin{proof}
(i) Each $N_k\sim N(0,1)$, so $P(|N_k|=c)=0$ and $P(\max_k|N_k|=c)=0$: $F_R$ is continuous. For $0\le c_1<c_2$, the law of $N$ is Gaussian with support the subspace $\operatorname{range}(R)\neq\{0\}$ (as $R_{kk}=1$), which contains a vector $x$ with $\max_k|x_k|=(c_1+c_2)/2$ by scaling. The set $\{c_1<\max_k|x_k|<c_2\}$ is open and meets the support, so it has positive probability, and $F_R(c_2)>F_R(c_1)$.
(ii) $F_R(c)\le P(|N_1|\le c)=2\Phi(c)-1$, which equals $1-\alpha$ at $c=z_{1-\alpha/2}$; so $F_R(z_{1-\alpha/2})\le1-\alpha=F_R(c_\alpha)$ and $z_{1-\alpha/2}\le c_\alpha$ by monotonicity. By the union bound, $1-F_R(c)\le\sum_kP(|N_k|>c)=2Q\{1-\Phi(c)\}$, which equals $\alpha$ at $c=z_{1-\alpha/(2Q)}$; so $F_R(z_{1-\alpha/(2Q)})\ge1-\alpha$.
(iii) $\{\max_{k\in S'}|N_k|\le c\}\subseteq\{\max_{k\in S}|N_k|\le c\}$ gives $F_{R_{S'}}\le F_{R_S}$, so $F_{R_S}(c_\alpha(R_{S'}))\ge1-\alpha=F_{R_S}(c_\alpha(R_S))$.
(iv) Let $R_n\to R$ and $c_n\to c$. With $\zeta\sim N(0,I_Q)$, $N^{(n)}=R_n^{1/2}\zeta\to R^{1/2}\zeta$ almost surely (the matrix square root is continuous), so $\max|N^{(n)}|-c_n\to\max|N|-c$ almost surely and, since $P(\max|N|=c)=0$, $F_{R_n}(c_n)\to F_R(c)$. Now let $c_n=c_\alpha(R_n)\in[z_{1-\alpha/2},z_{1-\alpha/(2Q)}]$; any limit point $c^\ast$ satisfies $F_R(c^\ast)=\lim F_{R_n}(c_n)=1-\alpha$, so $c^\ast=c_\alpha(R)$ by (i). Uniform continuity follows from compactness of $\mathcal R_Q$.
\end{proof}

\subsection{Proof of Theorem~\Mref{thm:fwer}}\label{si:testing-fwer}

Let $\mathcal Z=(Z_k(\theta_k))_{k\le Q}=(X_k/\varrho_k)_k$. By Theorem~\Mref{thm:joint} and $\varrho_k\to_p1$, along a subsequence with $R\to R_\infty$, $\mathcal Z\Rightarrow N(0,R_\infty)$, and $\hat R-R\to_p0$. By Lemma~\ref{lem:crit}(iv), $c_\alpha(\hat R_S)-c_\alpha(R_S)\to_p0$ for every fixed $S$ (a principal submatrix map is continuous), and $c_\alpha(R_S)\to c_\alpha(R_{\infty,S})$. Since limsup statements along the whole sequence are attained along subsequences, and there are finitely many sets $I_0$, it suffices to consider subsequences with $R\to R_\infty$ and a constant true set $I_0$. If $I_0=\emptyset$ there is nothing to prove.

\emph{(a), single step.} For $k\in I_0$, $\theta_k=\theta_{k,0}$ so $Z_k(\theta_{k,0})=\mathcal Z_k$. The false nulls do not enter $\mathcal Z_{I_0}$ or $\hat R$ (the critical value depends on the data only through $\hat R$, which is consistent for $R$ whatever the values of the $\theta_k$). Thus $(\max_{I_0}|\mathcal Z_k|,\,c_\alpha(\hat R))\Rightarrow(\max_{I_0}|N_k|,\,c_\alpha(R_\infty))$, the second component being constant in the limit and the first having a continuous law. By the portmanteau theorem,
\[
P\big(\max_{I_0}|\mathcal Z_k|>c_\alpha(\hat R)\big)\to P\big(\max_{I_0}|N_k|>c_\alpha(R_\infty)\big)\le P\big(\max_{I_0}|N_k|>c_\alpha(R_{\infty,I_0})\big)=\alpha,
\]
using Lemma~\ref{lem:crit}(iii) for the inequality. If $I_0=\{1,\dots,Q\}$ the inequality is an equality. This is the subset pivotality of the Gaussian limit: the limiting law of the vector of statistics of the true nulls does not depend on which other hypotheses are false.

\emph{(a), step-down.} The procedure sets $S_1=\{1,\dots,Q\}$, rejects $\{k\in S_j:|Z_k(\theta_{k,0})|>c_\alpha(\hat R_{S_j})\}$, sets $S_{j+1}=S_j$ minus the rejected, and stops when nothing is rejected. Let $j$ be the first step at which a true null is rejected. No true null was rejected before, so $I_0\subseteq S_j$, and $\max_{I_0}|\mathcal Z_k|>c_\alpha(\hat R_{S_j})\ge c_\alpha(\hat R_{I_0})$ by Lemma~\ref{lem:crit}(iii) applied to $\hat R$. Hence the familywise error event is contained in $\{\max_{I_0}|\mathcal Z_k|>c_\alpha(\hat R_{I_0})\}$, whose probability tends to $P(\max_{I_0}|N_k|>c_\alpha(R_{\infty,I_0}))=\alpha$ by the same portmanteau argument.

\emph{Bonferroni and Holm.} By Theorem~\Mref{thm:size}(a), for $k\in I_0$ and any fixed level $u$, $P(|\mathcal Z_k|>z_{1-u/2})\to u$. For Bonferroni, $P(\exists k\in I_0:|\mathcal Z_k|>z_{1-\alpha/(2Q)})\le\sum_{I_0}P(|\mathcal Z_k|>z_{1-\alpha/(2Q)})\to|I_0|\alpha/Q\le\alpha$. Holm's procedure in step-down form is that above with the critical value $z_{1-\alpha/(2|S_j|)}$ replacing $c_\alpha(\hat R_{S_j})$ (this rejects the same set as the one-at-a-time description with ordered $p$-values and thresholds $\alpha/(Q-j+1)$, because the thresholds are non-decreasing in $j$). If a true null is first rejected at step $j$ then $|S_j|\ge m_0=|I_0|$, so the critical value is at least $z_{1-\alpha/(2m_0)}$, and the probability of this event is at most $\sum_{I_0}P(|\mathcal Z_k|>z_{1-\alpha/(2m_0)})\to m_0\cdot\alpha/m_0=\alpha$.

\emph{(b).} The band covers all $\theta_k$ iff $\max_k|\mathcal Z_k|\le c_\alpha(\hat R)$, and
\[
P\big(\max_k|\mathcal Z_k|\le c_\alpha(\hat R)\big)\to F_{R_\infty}\big(c_\alpha(R_\infty)\big)=1-\alpha.
\]
 If every $\theta_k$ is covered then so is $\theta_{\hat k}$ for any data-dependent index $\hat k$.

\emph{(c).} The bounds are Lemma~\ref{lem:crit}(ii),(iii), valid for every $\hat R\in\mathcal R_Q$ and hence for every data set. Single-step max-$|t|$: $|Z_k|>z_{1-\alpha/(2Q)}\ge c_\alpha(\hat R)$ implies rejection. For the step-down procedures let $S^M_j,S^H_j$ be the active sets of max-$|t|$ and Holm. We show $S^M_j\subseteq S^H_j$ by induction; $S_1^M=S_1^H$. If it holds at $j$, then $c_\alpha(\hat R_{S^M_j})\le z_{1-\alpha/(2|S^M_j|)}\le z_{1-\alpha/(2|S^H_j|)}$ (using (ii) for the submatrix and $|S^M_j|\le|S^H_j|$). Take $k\in S^M_{j+1}$: then $k\in S^M_j\subseteq S^H_j$ and $|Z_k|\le c_\alpha(\hat R_{S^M_j})\le z_{1-\alpha/(2|S^H_j|)}$, so Holm does not reject $k$ at step $j$ and $k\in S^H_{j+1}$. Hence every hypothesis rejected by Holm is rejected by step-down max-$|t|$. Nested rejection regions give the power statement for every sample size. $\square$

\subsection{The population target: proof of Theorem~\Mref{thm:size}(d)}\label{si:testing-pop}

Fix one contrast. Write $\theta_U=f(T_U)=\phi(\bar t)$ for the pool contrast, with $\bar t=N^{-1}\sum_{i\le N}t_i$ the mean of the pool contributions, and $\theta_{\rm pop}=\phi(\tau)$, $\tau=E\,t_i$ (fixed, by (P1)). Put $\mathrm{er}_{\rm D}=\hat\theta-\theta_U$ (design error), $\mathrm{er}_{\rm P}=\theta_U-\theta_{\rm pop}$ (pool error), $V_D(u)=V_\psi(\theta;n)$ computed from the pool $u$, $\mathbf b=\nabla\phi(\tau)$, $\zeta_i=\mathbf b^\top t_i$ with $\operatorname{Var}\zeta_i=\sigma_\zeta^2>0$ by (P2), $v_{\rm P}=\sigma^2_\zeta/N$, $v_{\rm D}$ as in (P3) and $s^2=v_{\rm D}+v_{\rm P}$; $\varpi_N=v_{\rm P}/v_{\rm D}$.

\emph{Step 1: design-based statements, uniformly over $\mathcal E_N$.} Write $\mathcal E_N=\mathcal P_N$ for the class of (P2). For any sequence $u_N\in\mathcal P_N$, (A1)--(A4) hold (the constants $c,\varepsilon$ are common, $m_N\to\infty$ gives $\min n_h\to\infty$, $\eta_N\to0$ gives the dichotomy of (A4), and $\tau$ is handled along subsequences), so by Proposition~\Mref{prop:validity}(c),(d) (and (A5)), $\mathrm{er}_{\rm D}/V_D(u_N)^{1/2}\Rightarrow N(0,1)$ and $\hat V_D/V_D(u_N)\to_p1$ conditionally on the pool. By the argument of Corollary~\ref{cor:uniform} (this is where the quantitative class $\mathcal P_N$ replaces ``(A1)--(A4) hold'', which is a statement about sequences), for every $\eta>0$ and $t\in\mathbb R$,
\[
\sup_{u\in\mathcal E_N}\Big\{\sup_x\Big|P\Big(\tfrac{\mathrm{er}_{\rm D}}{V_D(u)^{1/2}}\le x\,\Big|\,u\Big)-\Phi(x)\Big|+P\Big(\Big|\tfrac{\hat V_D}{V_D(u)}-1\Big|>\eta\,\Big|\,u\Big)\Big\}\to0,
\]
and, by L\'evy's continuity theorem applied along nearly-worst sequences, for each $T$
\[
\sup_{u\in\mathcal E_N}\ \sup_{|\tau'|\le T}\Big|E\big(e^{i\tau'\mathrm{er}_{\rm D}/V_D(u)^{1/2}}\mid u\big)-e^{-\tau'^2/2}\Big|\to0.
\]

\emph{Step 2: the pool error.} By (P2) $\phi$ is $C^2$ with bounded derivatives on a neighbourhood of $\tau$ and $\|\bar t-\tau\|=O_p(N^{-1/2})$; Taylor's theorem gives $\mathrm{er}_{\rm P}=\mathbf b^\top(\bar t-\tau)+\Xi_N$ with $\Xi_N=O_p(N^{-1})$. By the classical central limit theorem for the i.i.d.\ bounded variables $\zeta_i$, $N^{1/2}\mathbf b^\top(\bar t-\tau)\Rightarrow N(0,\sigma^2_\zeta)$. Since $s^2\ge v_{\rm P}=\sigma_\zeta^2/N$, $\Xi_N/s=O_p(N^{-1/2})=o_p(1)$.

\emph{Step 3: the sum.} Fix $t\in\mathbb R$ and write $\varphi_N=E\exp\{it(\mathrm{er}_{\rm D}+\mathrm{er}_{\rm P})/s\}$. Condition on the pool:
\[
\varphi_N=E\Big[e^{it\,\mathrm{er}_{\rm P}/s}\,E\big(e^{it\,\mathrm{er}_{\rm D}/s}\mid u\big)\Big].
\]
Let $\Omega_N=\{u\in\mathcal E_N,\ |V_D(u)/v_{\rm D}-1|\le\tfrac12\}$; $P(\Omega_N)\to1$ by (P2) and (P3). On $\Omega_N$, $t\,V_D(u)^{1/2}/s\in[-2|t|,2|t|]$ (since $v_{\rm D}\le s^2$), so by Step 1, $E(e^{it\,\mathrm{er}_{\rm D}/s}\mid u)=\exp\{-t^2V_D(u)/(2s^2)\}+\epsilon_N(u)$ with $\sup_{u\in\Omega_N}|\epsilon_N(u)|\to0$. Hence
\[
\varphi_N=E\big[e^{it\,\mathrm{er}_{\rm P}/s}e^{-t^2V_D(u)/(2s^2)}\big]+o(1)=e^{-t^2v_{\rm D}/(2s^2)}E\,e^{it\,\mathrm{er}_{\rm P}/s}+o(1),
\]
where the last step uses $|e^{-t^2V_D/(2s^2)}-e^{-t^2v_{\rm D}/(2s^2)}|\le\{t^2/(2s^2)\}|V_D-v_{\rm D}|\wedge2\to_p0$ and bounded convergence (here (P3) is used: it is the only place where $V_D(u)$ is treated as non-random). By Step 2 and the uniform convergence of characteristic functions on compacts, with $m_N=1/(sN^{1/2})\le1/\sigma_\zeta$, $E\,e^{it\,\mathrm{er}_{\rm P}/s}=\exp\{-t^2m_N^2\sigma_\zeta^2/2\}+o(1)=\exp\{-t^2v_{\rm P}/(2s^2)\}+o(1)$. Therefore $\varphi_N=\exp\{-t^2(v_{\rm D}+v_{\rm P})/(2s^2)\}+o(1)=e^{-t^2/2}+o(1)$, that is, $(\hat\theta-\theta_{\rm pop})/s\Rightarrow N(0,1)$, with the remainder $\Xi_N/s$ already accounted for. No covariance between the two errors was computed; the conditional argument shows that the design error is asymptotically independent of the pool.

\emph{Step 4: variance estimators.} By Step 1 and (P3), $\hat V_D/v_{\rm D}\to_p1$. For the pool term, per-unit scale: $N\hat V_{\rm pool}=\hat g^\top\hat M\hat g$ where $\hat g=\nabla\phi(\hat T_\psi/N)$ and $\hat M$ is the weighted covariance matrix $N^{-1}\sum_{i\in s}w_it_it_i^\top-(N^{-1}\sum_{i\in s}w_it_i)(N^{-1}\sum_{i\in s}w_it_i)^\top$ of the labelled sample. Each entry of $N^{-1}\sum_{i\in s}w_it_{ic}t_{id}$ and of $N^{-1}\sum_{i\in s}w_it_{ic}$ is a stratified expansion estimator of the corresponding pool mean of a bounded variable, with conditional variance at most $Cr^2\to0$; the pool means of $t_{ic}t_{id}$ and $t_{ic}$ converge in probability to their population means by the law of large numbers. Hence $\hat M\to_p\operatorname{Cov}(t_i)$. Also $\hat T_\psi/N\to_p\tau$ (Proposition~\Mref{prop:validity}(b) and the law of large numbers), so $\hat g\to_pg$ and $N\hat V_{\rm pool}\to_pg^\top\operatorname{Cov}(t_i)g=\sigma^2_\zeta$, i.e.\ $\hat V_{\rm pool}/v_{\rm P}\to_p1$. Consequently $\varrho^2=(\hat V_D+\hat V_{\rm pool})/s^2$ lies between $\hat V_D/v_{\rm D}$ and $\hat V_{\rm pool}/v_{\rm P}$ and tends to one in probability.

\emph{Step 5.} Under $\theta_{\rm pop}=\theta_0$, $Z^{\rm pop}=\{(\hat\theta-\theta_0)/s\}/\varrho$ with $(\hat\theta-\theta_0)/s\Rightarrow N(0,1)$ and $\varrho\to_p1$, which are the two inputs of Lemma~\ref{lem:slutsky}; the proofs of Theorem~\Mref{thm:size}(a)--(c) and Theorem~\Mref{thm:power} then apply verbatim to $Z^{\rm pop}$ (with $\mu=(\theta_{\rm pop}-\theta_0)/s$), which gives in particular the population power statements of Theorem~\Mref{thm:power}; $\alpha<1/2$ is needed for the one-sided and TOST statements because Lemma~\ref{lem:slutsky} requires $z>0$. For the pool-conditional statistic, $Z=\{(\hat\theta-\theta_0)/s\}\cdot s/\hat V_D^{1/2}$ and $s^2/\hat V_D\to_p1+\varpi_N$. Along a subsequence with $\varpi_N\to\varpi_\infty<\infty$, $Z\Rightarrow N(0,1+\varpi_\infty)$ and the rejection probability tends to $P(|N(0,1)|>z_{1-\alpha/2}/\sqrt{1+\varpi_\infty})=2\{1-\Phi(z_{1-\alpha/2}/\sqrt{1+\varpi_\infty})\}$. If $\varpi_N\to\infty$ then $|Z|\ge|(\hat\theta-\theta_0)/s|\cdot s/\hat V_D^{1/2}$ with the first factor converging to $|N(0,1)|$, which has no atom at zero, and the second tending to infinity in probability, so $P(|Z|>z)\to1$, consistent with the formula at $\varpi_\infty=\infty$. $\square$

\subsection{Proof of Theorem~\Mref{thm:design}}\label{si:testing-design}

Write $F(n,w)=\sum_kw_kV_k(n)/\delta_k^2$, $\mathcal N=\{n:\sum_hn_h=L,\ \ell_h\le n_h\le N_h\}$ (compact, convex, non-empty by $\sum\ell_h\le L\le\sum N_h$), and note $V_k(n)=\sum_ha_{hk}/n_h-b_k$ with $a_{hk}=N_h^2S^2_{hk}\ge0$, $b_k=\sum_hN_hS^2_{hk}$, so $V_k$ is convex, and $F(n,w)$ equals the objective $\Phi(n)$ of Proposition~\Mref{prop:allocation} with weights $\omega=w/\delta^2$ (we write $F(\cdot,w)$ for it; the symbol $\Phi$ without argument $n$ is the normal distribution function).

\emph{(a)} By Theorem~\Mref{thm:power} the asymptotic power depends on the design only through $\mu=\Upsilon/V^{1/2}$ and $\Psi_\alpha$ is strictly increasing in $|\mu|$ (Supplementary Section~\SIref{si:testing-are}); the one-sided power $\Phi(\mu-z)$ is increasing in $\mu$. So any $n$ minimising $V$ maximises power at every $\Upsilon$. With $Q=1$, Proposition~\Mref{prop:allocation}(b) gives the minimiser $n(\omega)$ (weights $\omega=1/\delta^2$ do not change the allocation).

\emph{(b)} \emph{Saddle point.} $F$ is continuous, convex in $n$ and linear in $w$, and both $\mathcal N$ and the simplex $\mathcal S_Q$ are compact and convex. Sion's minimax theorem \citep{sion1958} gives $\min_n\max_wF=\max_w\min_nF$. Now $\max_{w\in\mathcal S_Q}F(n,w)=\max_kV_k(n)/\delta_k^2$, a continuous function with a minimiser $n^\circ$, and $\Gamma(w)=\min_nF(n,w)$ is a minimum of affine functions, hence concave and upper semi-continuous, with a maximiser $w^\circ$ on $\mathcal S_Q$. Let $t^\ast=\Gamma(w^\circ)=\max_kV_k(n^\circ)/\delta_k^2$. Then $t^\ast=\Gamma(w^\circ)\le F(n^\circ,w^\circ)\le\max_kV_k(n^\circ)/\delta_k^2=t^\ast$, so equality holds throughout: $n^\circ\in\arg\min_nF(\cdot,w^\circ)$, and $F(n^\circ,w^\circ)=\max_wF(n^\circ,w)$, which means that $w^\circ$ is supported on $\{k:V_k(n^\circ)/\delta_k^2=t^\ast\}$ (complementary slackness). Conversely, any pair $(n,w)$ with $n\in\arg\min F(\cdot,w)$ and $w$ supported on the argmax set of $V_k(n)/\delta_k^2$ is optimal: for every $n'$, $\max_kV_k(n')/\delta_k^2\ge F(n',w)\ge F(n,w)=\max_kV_k(n)/\delta_k^2$.
\emph{The primal solution from the dual one.} The argument above shows that for every maximiser $w^\ast$ of $\Gamma$ and every primal optimum $n^\circ$, $n^\circ\in\arg\min F(\cdot,w^\ast)$ (the equality chain with $w^\circ$ replaced by $w^\ast$, which has $\Gamma(w^\ast)=t^\ast$). Now $\arg\min F(\cdot,w^\ast)=\arg\min\Phi_{\omega^\ast}$; if $A_h(\omega^\ast)>0$ for all $h$, $\Phi_{\omega^\ast}(n)=\sum_hN_h^2A_h(\omega^\ast)^2/n_h-\text{const}$ is strictly convex, so the minimiser is unique and equals $n(\omega^\ast)$ by Proposition~\Mref{prop:allocation}(b). Thus the optimum is unique and given by (\Mref{eq:n-omega}). Without bounds, Proposition~\Mref{prop:allocation}(a) gives $\Gamma(w)=(\sum_hN_hA_h)^2/L-\sum_hN_hA_h^2$ with $A_h^2=\sum_kw_kS^2_{hk}/\delta_k^2$.
\emph{Gradient.} For $w$ with $A_h(w/\delta^2)>0$ for all $h$ the minimiser $n(w/\delta^2)$ of $F(\cdot,w)$ over the compact $\mathcal N$ is unique; the superdifferential of the concave function $\Gamma$ is the closed convex hull of the gradients $\nabla_wF(n,w)=(V_k(n)/\delta_k^2)_k$ over minimisers $n$ (Danskin's theorem), here the single vector $(V_k(n(w/\delta^2))/\delta_k^2)_k$. So $\Gamma$ is differentiable there with this gradient, and the first-order condition for maximising $\Gamma$ on $\mathcal S_Q$ ($w_k>0\Rightarrow\partial_kg=\max_j\partial_jg$) is the complementary slackness above.
\emph{Support bound.} Let $(n^\ast,w^\ast)$ be a saddle point and $J=\operatorname{supp}w^\ast$. Since $n^\ast_h\ge\ell_h>0$, $\Phi_{\omega}$ is differentiable at $n^\ast$ with $\partial_h\Phi_\omega=-\sum_k(w_k/\delta_k^2)N_h^2S^2_{hk}/(n^\ast_h)^2=-\sum_kw_kG_{hk}$, $G_{hk}=N_h^2S^2_{hk}/\{\delta_k^2(n^\ast_h)^2\}$. Because $\Phi_{\omega^\ast}$ is convex, $n^\ast$ minimises it over the polytope $\mathcal N$ iff $\sum_h(\sum_kw^\ast_kG_{hk})(n_h-n^\ast_h)\le0$ for all $n\in\mathcal N$. Every $n-n^\ast$ has coordinate sum zero, so the condition depends on the vector $p=\sum_{k\in J}w^\ast_kG_{\cdot k}\in\mathbb R^H$ only modulo $\mathbb R\mathbf 1$. In the quotient $\mathbb R^H/\mathbb R\mathbf1\cong\mathbb R^{H-1}$ the image of $p$ lies in the convex hull of the images of $\{G_{\cdot k}:k\in J\}$, so by Carath\'eodory's theorem it is a convex combination, with weights $w'$, of at most $H$ of them. Then $w'\in\mathcal S_Q$ is supported in $J$, $\sum_kw'_kG_{\cdot k}-p\in\mathbb R\mathbf1$, hence $n^\ast$ also minimises $F(\cdot,w')$, and $w'$ is supported on the argmax set, so $(n^\ast,w')$ is a saddle point and $w'$ maximises $\Gamma$.

\emph{(c)} Let $n^{\rm def}=n(\omega^{\rm def})$ be built from the same array $S^2_{hk}$ that defines the $V_k$. By Proposition~\Mref{prop:allocation}(b) it minimises $\sum_kV_k(n)/\sigma_k^2$ over $\mathcal N$; with $V^{\rm SRS}_k=N^2\sigma_k^2(1+O(1/N))/L$ (up to the finite-population correction, common to all $k$) this is the sum of inverse efficiencies up to a factor not depending on $k$ or $n$ (Proposition~\Mref{prop:allocation}(d)). If the paper's allocation (computed from plain variances) is used with the difference estimator, $V_k$ is the residual variance and that allocation is not $n(\omega^{\rm def})$ for the residual array; none of the statements of (c) then applies to it, and only the weighted-sum loss bound of Proposition~\SIref{prop:alloc-dif} is available. \emph{Sufficiency.} Suppose $V_k(n^{\rm def})/\delta_k^2=t$ for all $k$. Put $w_k\propto\delta_k^2\omega^{\rm def}_k$ on the simplex. The allocation (\Mref{eq:n-omega}) is invariant under a common rescaling of $\omega$ ($\lambda$ rescales accordingly), so $n(w/\delta^2)=n^{\rm def}$, and every $k$ is in the argmax set; by the converse in (b), $n^{\rm def}$ is optimal. \emph{Necessity.} Assume $n^{\rm def}$ is optimal, no stratum is at a bound, $A_h(\omega^{\rm def})>0$ for all $h$, and the vectors $(S^2_{hk})_{h\le H}$, $k\le Q$, are linearly independent. By (b) there is a maximiser $w^\ast$ with $n^{\rm def}=n(\omega^\ast)$ (we also assume $A_h(\omega^\ast)>0$ for all $h$, so that the minimiser is unique). No stratum being at a bound, $n_h\propto N_hA_h$, so $A_h(\omega^\ast)^2=cA_h(\omega^{\rm def})^2$ for all $h$, i.e.\ $\sum_k(\omega^\ast_k-c\,\omega^{\rm def}_k)S^2_{hk}=0$ for all $h$; linear independence of the $Q$ vectors forces $\omega^\ast=c\,\omega^{\rm def}$, so $w^\ast_k\propto\delta_k^2/\sigma_k^2>0$ for every $k$, and complementary slackness makes all $V_k(n^{\rm def})/\delta_k^2$ equal. (Linear independence needs $Q\le H$. For $Q>H$ equal relative variances are not necessary: a numerical example with $H=2$, $Q=3$, in which $n^{\rm def}$ has relative variances $(1,1,0.01)$ and is maximin-optimal, is in \texttt{docs/testing\_proof\_check.md}; it is not proved here. For the full family, optimality of $n^{\rm def}$ has to be checked numerically from the saddle-point conditions.) \emph{Certificate.} For any $w\in\mathcal S_Q$ and any $n$, weak duality gives $\Gamma(w)\le\min_{n'}\max_kV_k(n')/\delta_k^2=t^\ast\le\max_kV_k(n)/\delta_k^2$, so the ratio $\max_kV_k(n^{\rm def})/\delta_k^2$ to $\Gamma(w)$ bounds the loss of the default (the factor by which the labels needed for a common noncentrality exceed the optimum, when finite-population corrections are negligible).

\emph{(d)} For $\tilde n$ with $\tilde n_h\ge n_h-1$ and $n_h\ge n_0>1$, $V_k(\tilde n)+b_k=\sum_ha_{hk}/\tilde n_h\le\sum_ha_{hk}/(n_h-1)=\sum_h\{n_h/(n_h-1)\}a_{hk}/n_h\le\{n_0/(n_0-1)\}\sum_ha_{hk}/n_h=\{n_0/(n_0-1)\}\{V_k(n)+b_k\}$, because $x\mapsto x/(x-1)$ is decreasing on $(1,\infty)$. An integer allocation with $\tilde n_h\in\{\lfloor n_h\rfloor,\lceil n_h\rceil\}$ and $\sum_h\tilde n_h=L$ exists, satisfies $\tilde n_h\ge n_h-1$, and $\tilde n_h\le\lceil n_h\rceil\le N_h$ and $\tilde n_h\ge\ell_h$ if the bounds are integers. $\square$

\subsection{Proof of Proposition~\Mref{prop:eo}}\label{si:testing-eo}

Let $(\hat a,\hat b)$ be the plug-in gaps with variances $V_a,V_b\asymp r^2$ (assumption (A3) for the two gaps, which are smooth functions of the totals like the comparison contrasts) and suppose the standardised vector $((\hat a-a)/V_a^{1/2},(\hat b-b)/V_b^{1/2})$ is asymptotically $N(0,R_\infty)$ along subsequences (Theorem~\Mref{thm:joint} with $Q=2$), so that $Z^N_a=(\hat a-a)/r$ and $Z_b^N=(\hat b-b)/r$ are tight with a Gaussian limit $(Z_a,Z_b)$ along subsequences. Let $s_a=\operatorname{sgn}a$, $s_b=\operatorname{sgn}b$.

\emph{(a)} Take $|a|>|b|$ (the other case is symmetric) and put $\varsigma_N=|a|-|b|$, with $\varsigma_N/r\to\infty$. On $\Omega_N=\{|\hat a-a|<\varsigma_N/4,\ |\hat b-b|<\varsigma_N/4\}$, whose probability tends to one because $|\hat a-a|+|\hat b-b|=O_p(r)=o_p(\varsigma_N)$, we have $|\hat a-a|<\varsigma_N\le|a|$, so $\operatorname{sgn}\hat a=s_a$, and $|\hat a|\ge|a|-\varsigma_N/4>|b|+\varsigma_N/4>|\hat b|$. Hence on $\Omega_N$, $\widehat{\mathrm{EO}}=|\hat a|=s_a\hat a$ and $\widehat{\mathrm{EO}}-\mathrm{EO}=s_a(\hat a-a)$. Therefore $(\widehat{\mathrm{EO}}-\mathrm{EO})/V_a^{1/2}=s_a(\hat a-a)/V_a^{1/2}+o_p(1)\Rightarrow N(0,1)$ by Theorem~\Mref{thm:joint} (symmetry of $N(0,1)$); in the case $|b|>|a|$ the same argument gives the limit with $V_b$. With $\hat V_a/V_a\to_p1$ (Proposition~\Mref{prop:validity}(d)) the Wald interval has asymptotic coverage $1-\alpha$ by Slutsky's lemma.

\emph{(b)} Let $\vartheta_N=(|a|-|b|)/r\to\vartheta\in\mathbb R$ (either sign) and $|a|/r,|b|/r\to\infty$. Then $P(\operatorname{sgn}\hat a=s_a,\operatorname{sgn}\hat b=s_b)\to1$, since $|\hat a-a|=O_p(r)=o_p(|a|)$ and likewise for $b$. On this event $|\hat a|=|a|+s_a(\hat a-a)$ and $|\hat b|=|b|+s_b(\hat b-b)=|a|-r\vartheta_N+s_b(\hat b-b)$, so $(\widehat{\mathrm{EO}}-|a|)/r=\max(s_aZ_a^N,s_bZ_b^N-\vartheta_N)$, and $\mathrm{EO}=\max(|a|,|b|)=|a|+r\max(0,-\vartheta_N)$. Therefore
\[
\frac{\widehat{\mathrm{EO}}-\mathrm{EO}}{r}=\max\big(s_aZ_a^N,\ s_bZ_b^N-\vartheta_N\big)-\max(0,-\vartheta_N)\Rightarrow\max(s_aZ_a,\ s_bZ_b-\vartheta)-\max(0,-\vartheta)
\]
by the continuous mapping theorem. With $U_1=s_aZ_a$, $U_2=s_bZ_b-\vartheta$ we have $\max(U_1,U_2)=(U_1+U_2)/2+|U_1-U_2|/2$, $EU_1=0$, $EU_2=-\vartheta$, and $U_1-U_2=Y\sim N(\vartheta,\sigma^2)$, $\sigma^2=\operatorname{Var}(s_aZ_a-s_bZ_b)$. Hence the mean of the limit is $-\vartheta/2+E|Y|/2-\max(0,-\vartheta)=\tfrac12\{E|Y|-|\vartheta|\}$ (check $\vartheta\ge0$ and $\vartheta<0$ separately), and $E|Y|\ge|EY|=|\vartheta|$ with strict inequality when $\sigma>0$. (A Monte Carlo check with correlated Gaussian errors and $\vartheta\in\{-2,-1,0,0.7,1\}$ agrees with this mean to three decimals. The bound $\max(0,-\vartheta)$ of an earlier version referred to $(\widehat{\mathrm{EO}}-|a|)/r$, not to $(\widehat{\mathrm{EO}}-\mathrm{EO})/r$.) The limit is not normal in general: for $\vartheta=0$ and independent standard normal $s_aZ_a,s_bZ_b$ the maximum has density $2\varphi(x)\Phi(x)$ and mean $\pi^{-1/2}>0$. At $a=b=0$, $\widehat{\mathrm{EO}}=\max(|\hat a|,|\hat b|)=r\max(|Z_a^N|,|Z_b^N|)$ exactly, which gives the stated limit. The regime with $|a|/r$ or $|b|/r$ bounded but not both zero needs the limit of $\max(|a+rZ_a|,|b+rZ_b|)$ with finite offsets and is not treated.

\emph{(c)} $\{\mathrm{EO}_m\ge\delta\}=\{a\ge\delta\}\cup\{a\le-\delta\}\cup\{b\ge\delta\}\cup\{b\le-\delta\}$. The procedure rejects iff the four one-sided tests reject: $Z_a(-\delta)>z_{1-\alpha}$, $Z_a(\delta)<-z_{1-\alpha}$, and the same for $b$; equivalently the $(1-2\alpha)$ Wald intervals of $a$ and $b$ lie in $(-\delta,\delta)$. If $\mathrm{EO}_m\ge\delta$ at least one of the four null hypotheses holds; pass to a subsequence along which it is always the same one. By Theorem~\Mref{thm:size}(c) (and its mirror image for the lower margin) the corresponding one-sided test has $\limsup P(\text{reject})\le\alpha$, and the intersection test rejects only if this test does. The argument uses only the marginal assumptions for $a$ and $b$ and no condition on ties. $\square$
